\documentclass[tikz,border=4pt]{standalone}
\usepackage{ctex}
\usepackage{amsmath,amssymb}
\usepackage{pgfplots}
\pgfplotsset{compat=1.18}
\usepgfplotslibrary{groupplots,fillbetween}
\begin{document}
\begin{tikzpicture}
\begin{groupplot}[
  group style={
    group size=2 by 1,
    horizontal sep=2.0cm,
  },
  width=6.5cm, height=5.5cm,
  xlabel={$x$},
  xlabel style={font=\small},
  ylabel style={font=\small},
  tick label style={font=\footnotesize},
  grid=both, grid style={gray!20},
  axis lines=left,
  samples=200,
  legend style={font=\footnotesize, at={(0.98,0.98)}, anchor=north east},
]

% Plot 1: Two distributions p and q
\nextgroupplot[
  title={分布 $p$ 与近似分布 $q$},
  title style={font=\small},
  xmin=-4, xmax=6,
  ymin=-0.02, ymax=0.55,
  ylabel={密度},
]
% p: narrow Gaussian at mu=1
\addplot[name path=p, blue, very thick, domain=-4:6]
  {exp(-(x-1)^2/0.5)/sqrt(0.5*3.14159)};
\addlegendentry{$p(x)=\mathcal{N}(1,0.25)$（真实）}
% q: wider Gaussian at mu=1.5
\addplot[name path=q, red, very thick, domain=-4:6]
  {exp(-(x-1.5)^2/2.0)/sqrt(2.0*3.14159)};
\addlegendentry{$q(x)=\mathcal{N}(1.5,1)$（近似）}
% fill difference region
\addplot[blue!15, opacity=0.4] fill between [of=p and q, soft clip={domain=-1:4}];
% KL formula
\node[font=\footnotesize, fill=white, draw=blue!50, rounded corners, align=center]
  at (axis cs:4.0, 0.4)
  {$\mathrm{KL}(p\|q)=\int p\ln\frac{p}{q}\,dx$\\[2pt]$\geq 0$，等号 $\Leftrightarrow$ $p=q$};

% Plot 2: KL integrand p*log(p/q)
\nextgroupplot[
  title={KL 被积函数 $p(x)\ln\frac{p(x)}{q(x)}$},
  title style={font=\small},
  xmin=-3, xmax=5,
  ymin=-0.1, ymax=0.5,
  ylabel={$p\ln(p/q)$},
]
% KL integrand: compute p*log(p/q) where p=N(1,0.25), q=N(1.5,1)
% log(p/q) = log(p) - log(q)
% = -0.5*(x-1)^2/0.25 - 0.5*log(0.5*pi) + 0.5*(x-1.5)^2/2 + 0.5*log(2*pi)
% = -2(x-1)^2 + 0.25(x-1.5)^2 + 0.5*log(4)
% p*log(p/q) = p * above
\addplot[name path=kld, purple, very thick, domain=-3:5]
  { exp(-(x-1)^2/0.5)/sqrt(0.5*3.14159) *
    (-2*(x-1)^2 + 0.25*(x-1.5)^2 + 0.6931) };
\addlegendentry{$p(x)\ln\frac{p(x)}{q(x)}$}
\addplot[name path=zero3, draw=none, domain=-3:5] {0};
\addplot[purple!15] fill between [of=kld and zero3];
% area = KL value
\node[font=\footnotesize, fill=white, draw=purple!40, rounded corners, align=center]
  at (axis cs:3.5, 0.35)
  {面积 $=\mathrm{KL}(p\|q)$\\（非负！）};
% asymmetry note
\node[font=\footnotesize, draw=orange!70, fill=orange!10, rounded corners, align=center]
  at (axis cs:-1.0, 0.35)
  {注意：\\$\mathrm{KL}(p\|q)\neq\mathrm{KL}(q\|p)$\\KL 非对称！};

\end{groupplot}

\node[font=\normalsize\bfseries] at (8.5, 6.3) {KL 散度：两个分布差异的积分度量};
\end{tikzpicture}
\end{document}
